% ==============================================================================
% COMMANDS: THÀNH PHẦN NHANH
%  (Xem Phụ lục B của bản PDF để có ví dụ + mã nguồn của từng thành phần)
% ==============================================================================
\makeatletter

% ==============================================================================
% 1. NHẤN MẠNH CHỮ TRONG ĐOẠN
% ==============================================================================
\newcommand{\nhanmanh}[1]{\textbf{\textcolor{bkxanh}{#1}}}         % ý chính
\newcommand{\quantrong}[1]{\textbf{\textcolor{do}{#1}}}            % rất quan trọng
\newcommand{\danhdau}[1]{\colorbox{vangnhat}{\strut #1}}           % bút dạ quang (cụm ngắn)
% Thuật ngữ: \thuatngu{Trí tuệ nhân tạo}[Artificial Intelligence]
\NewDocumentCommand{\thuatngu}{m o}{%
  \textbf{\textcolor{bkxanhdam}{#1}}\IfValueT{#2}{ (\textit{#2})}}
% Nhãn nhỏ: \nhan{Mới}  \nhan[do]{Quan trọng}
\NewDocumentCommand{\nhan}{O{bkxanhsang} m}{%
  \tcbox[on line, colback=#1, colframe=#1, colupper=white, boxrule=0pt,
         arc=2pt, boxsep=0pt, left=3pt, right=3pt, top=1.5pt, bottom=1.5pt,
         fontupper=\footnotesize\bfseries]{#2}}
% Ghi chú "cần bổ sung" để nhắc việc khi viết nháp
\newcommand{\canbosung}[1]{\nhan[do]{CẦN BỔ SUNG}~\textcolor{do}{\textit{#1}}}
% Biểu tượng tròn: \bieutuong{xanhla}{\ding{51}}
\newcommand{\bieutuong}[2]{%
  \tikz[baseline=-0.6ex]\node[circle, fill=#1, text=white, inner sep=0pt,
    minimum size=1.15em, font=\footnotesize\bfseries]{#2};}
\newcommand{\dung}{\textcolor{xanhla}{\ding{51}}}
\newcommand{\sai}{\textcolor{do}{\ding{55}}}

% Danh sách có dấu tick / dấu chéo / ngôi sao
\newlist{dsdung}{itemize}{1}\setlist[dsdung]{label=\dung, leftmargin=1.8em}
\newlist{dssai}{itemize}{1}\setlist[dssai]{label=\sai, leftmargin=1.8em}
\newlist{dssao}{itemize}{1}\setlist[dssao]{label=\textcolor{cam}{\ding{72}}, leftmargin=1.8em}

% ==============================================================================
% 2. HỘP / KHUNG CHỮ (tcolorbox)
% ==============================================================================
% Kiểu chung: vạch màu bên trái + dải tiêu đề nhạt
\tcbset{
  hopvach/.style={enhanced, breakable, colback=#1nhat, colframe=#1,
    boxrule=0pt, leftrule=3.5pt, arc=0pt, outer arc=0pt,
    left=9pt, right=9pt, top=5pt, bottom=6pt, before skip=8pt, after skip=8pt,
    colbacktitle=#1!16!white, coltitle=#1dam, fonttitle=\bfseries,
    titlerule=0pt, toptitle=2pt, bottomtitle=2pt},
  hopdac/.style={enhanced, breakable, colback=white, colframe=#1,
    colbacktitle=#1, coltitle=white, fonttitle=\bfseries, boxrule=0.8pt,
    arc=1.5mm, left=9pt, right=9pt, top=6pt, bottom=6pt,
    before skip=8pt, after skip=8pt},
}
% 2.1 Hộp thông báo  — \begin{ghichu}[Tiêu đề tùy chọn] ... \end{ghichu}
\newtcolorbox{ghichu}[1][Ghi chú]{hopvach=xanhduong,
  title={\bieutuong{xanhduong}{i}\hspace{5pt}#1}}
\newtcolorbox{luuy}[1][Lưu ý]{hopvach=cam,
  title={\bieutuong{cam}{!}\hspace{5pt}#1}}
\newtcolorbox{canhbao}[1][Cảnh báo]{hopvach=do,
  title={\bieutuong{do}{\ding{55}}\hspace{5pt}#1}}
\newtcolorbox{goiy}[1][Gợi ý]{hopvach=xanhla,
  title={\bieutuong{xanhla}{\ding{51}}\hspace{5pt}#1}}

% 2.2 Hộp kết quả / kết luận nổi bật
\newtcolorbox{ketqua}[1][Kết quả chính]{hopdac=xanhla, colback=xanhlanhat, title={#1}}

% 2.3 Hộp đánh số theo chương: Định nghĩa 2.1, Định lý 2.1, Ví dụ 2.1
%     \begin{dinhnghia}[Tên tùy chọn] ... \label{dn:x} \end{dinhnghia}
\newtcolorbox[auto counter, number within=chapter]{dinhnghia}[1][]{hopdac=bkxanh,
  colback=bkxanhnen,
  title={Định nghĩa~\thetcbcounter\ifstrempty{#1}{}{:~#1}}}
\newtcolorbox[auto counter, number within=chapter]{dinhly}[1][]{hopdac=tim,
  colback=timnhat,
  title={Định lý~\thetcbcounter\ifstrempty{#1}{}{:~#1}}}
\newtcolorbox[auto counter, number within=chapter]{vidu}[1][]{hopvach=ngoc,
  title={Ví dụ~\thetcbcounter\ifstrempty{#1}{}{:~#1}}}

% 2.4 Khung chữ tự chọn màu: \begin{khungchu}[màu]{Tiêu đề}  hoặc  [màu] không tiêu đề
\NewTColorBox{khungchu}{O{bkxanh} o}{hopdac=#1, colback=#1nhat,
  IfValueT={#2}{title={#2}}}

% 2.5 Khung viền nét đứt (đề bài, yêu cầu, trích đoạn)
\NewTColorBox{khungnetdut}{O{bkxanh}}{enhanced, breakable, colback=white,
  frame hidden, borderline={0.8pt}{0pt}{#1, dashed}, arc=1.5mm,
  left=9pt, right=9pt, top=6pt, bottom=6pt}

% 2.5b Lời giải (khung trắng xám, chữ xanh dương):  \begin{loigiai}[Tiêu đề] ... \end{loigiai}
\newtcolorbox{loigiai}[1][Lời giải]{enhanced, breakable,
  colback=xam!4, colframe=xam!35, boxrule=0.6pt, arc=1.5mm,
  colupper=xanhduongdam, left=10pt, right=10pt, top=6pt, bottom=6pt,
  before skip=8pt, after skip=8pt,
  colbacktitle=xam!12, coltitle=xanhduongdam, fonttitle=\bfseries, titlerule=0.4pt,
  title={\ding{46}\hspace{5pt}#1},
  after upper={\par\nobreak\hfill\textcolor{xam!70}{$\blacksquare$}}}

% 2.6 Trích dẫn: \begin{trichdan}{Tác giả, năm} ... \end{trichdan}
\newtcolorbox{trichdan}[1]{enhanced, breakable, colback=xamnhat, frame hidden,
  boxrule=0pt, arc=0pt, left=26pt, right=12pt, top=8pt, bottom=8pt,
  fontupper=\itshape,
  overlay={\node[anchor=north west, text=bkxanhsang!60, font=\fontsize{40}{40}\selectfont]
           at ([xshift=2pt, yshift=6pt]frame.north west) {``};
           \draw[bkxanhsang, line width=2.5pt] (frame.north west) -- (frame.south west);},
  after upper={\par\smallskip\hfill\normalfont\small --- #1}}

% 2.7 Tóm tắt / ý chính của chương
\newtcolorbox{ychinh}[1][Ý chính của chương]{hopdac=bkxanh, colback=bkxanhnen,
  title={\ding{72}\hspace{5pt}#1}}

% 2.8 So sánh hai cột: \SoSanh[màu1][màu2]{Tiêu đề 1}{Nội dung 1}{Tiêu đề 2}{Nội dung 2}
\NewDocumentCommand{\SoSanh}{O{xanhla} O{do} m +m m +m}{%
  \begin{tcbraster}[raster columns=2, raster equal height, raster column skip=8pt,
                    raster before skip=8pt, raster after skip=8pt]
    \begin{tcolorbox}[hopdac=#1, unbreakable, colback=#1nhat, title={#3}]#4\end{tcolorbox}%
    \begin{tcolorbox}[hopdac=#2, unbreakable, colback=#2nhat, title={#5}]#6\end{tcolorbox}%
  \end{tcbraster}}

% 2.9 Ô chỉ số (số liệu nổi bật): \begin{hangchiso}[3] \ChiSo{16\%}{Mô tả} ... \end{hangchiso}
\newenvironment{hangchiso}[1][3]
  {\begin{tcbraster}[raster columns=#1, raster equal height, raster column skip=8pt,
                     raster before skip=8pt, raster after skip=8pt]}
  {\end{tcbraster}}
\NewDocumentCommand{\ChiSo}{O{bkxanh} m m}{%
  \begin{tcolorbox}[enhanced, colback=#1nhat, colframe=#1, boxrule=0pt,
      toprule=3pt, arc=1mm, halign=center, valign=center, top=6pt, bottom=6pt]
    {\fontsize{22}{26}\selectfont\bfseries\color{#1}#2}\par\smallskip
    {\small #3}
  \end{tcolorbox}}

% ==============================================================================
% 3. CÔNG THỨC
% ==============================================================================
% 3.1 Khung công thức có tiêu đề: \begin{khungcongthuc}[Tiêu đề] \begin{equation}...\end{equation} \end{khungcongthuc}
\NewTColorBox{khungcongthuc}{o}{enhanced, colback=vangnhat!60!white,
  colframe=cam, boxrule=0.6pt, arc=1.5mm, left=6pt, right=6pt, top=2pt, bottom=4pt,
  before skip=8pt, after skip=8pt,
  IfValueT={#1}{title={#1}, colbacktitle=cam, coltitle=white, fonttitle=\bfseries\small}}
% 3.2 Đóng khung một công thức ngay trong equation:  \begin{equation} \NoiBat{E=mc^2} \end{equation}
\tcbset{highlight math style={enhanced, colback=vangnhat, colframe=cam,
  boxrule=0.6pt, arc=1mm, left=4pt, right=4pt, top=3pt, bottom=3pt}}
\newcommand{\NoiBat}[1]{\tcbhighmath{#1}}
% 3.3 Ký hiệu thường dùng
\newcommand{\R}{\mathbb{R}}
\newcommand{\N}{\mathbb{N}}
\newcommand{\Z}{\mathbb{Z}}
\newcommand{\dif}{\mathop{}\!\mathrm{d}}
\DeclarePairedDelimiter{\abs}{\lvert}{\rvert}
\DeclarePairedDelimiter{\norm}{\lVert}{\rVert}

% ==============================================================================
% 4. BẢNG BIỂU
% ==============================================================================
% Kiểu cột: L/C/R{rộng} canh trái/giữa/phải;  Y = tabularx giãn, canh trái;  Z = giãn, canh giữa
\newcolumntype{L}[1]{>{\raggedright\arraybackslash}p{#1}}
\newcolumntype{C}[1]{>{\centering\arraybackslash}p{#1}}
\newcolumntype{R}[1]{>{\raggedleft\arraybackslash}p{#1}}
\newcolumntype{Y}{>{\raggedright\arraybackslash}X}
\newcolumntype{Z}{>{\centering\arraybackslash}X}
% Hàng tiêu đề: \hangtieude \thead{Cột 1} & \thead{Cột 2} \\
\renewcommand{\theadfont}{\bfseries\color{white}}
\renewcommand{\theadalign}{cc}
\newcommand{\hangtieude}{\rowcolor{bkxanh}}
% \thead dựng một bảng con bên trong ô; khi bảng ngoài dùng \rowcolors (bangsoc,
% danh mục...) bảng con bị tô màu sọc và làm lệch thứ tự sọc -> tạm tắt rồi khôi phục.
\let\thead@goc\thead
\renewcommand{\thead}{\@ifnextchar[\thead@moi{\thead@moi[]}}
\def\thead@moi[#1]#2{%
  \let\thead@rc\if@rowcolors \edef\thead@rn{\the\rownum}%
  \global\@rowcolorsfalse
  \ifx\relax#1\relax\thead@goc{#2}\else\thead@goc[#1]{#2}\fi
  \global\rownum=\thead@rn\relax
  \ifx\thead@rc\iftrue\global\@rowcolorstrue\fi}
% Ô nổi bật trong bảng
\newcommand{\otot}{\cellcolor{xanhlanhat}}
\newcommand{\oxau}{\cellcolor{donhat}}
\newcommand{\onoibat}{\cellcolor{vangnhat}}
% Bảng nhanh:  \begin{bangbieu}{Chú thích}{tab:nhan}{C{1cm} Y Y} ... \end{bangbieu}
\newenvironment{bangbieu}[3]
  {\table[H]\centering\caption{#1}\label{#2}%
   \tabularx{\linewidth}{#3}\toprule}
  {\bottomrule\endtabularx\endtable}
% Bảng sọc (hàng chẵn/lẻ xen màu)
\newenvironment{bangsoc}[3]
  {\table[H]\centering\caption{#1}\label{#2}%
   \rowcolors{2}{white}{bkxanhnen}\tabularx{\linewidth}{#3}\toprule}
  {\bottomrule\endtabularx\endtable}
% Ghi chú/nguồn dưới bảng: \nguonbang{Tổng hợp từ [1]}
\newcommand{\nguonbang}[1]{\par\vspace{-4pt}\noindent{\footnotesize\textit{Nguồn:} #1}\par}

% ==============================================================================
% 5. HÌNH ẢNH
% ==============================================================================
% Khung ảnh: viền xám mảnh, đệm trắng, bo góc
\tcbset{khunganh/.style={enhanced, colback=white, colframe=xam!35, boxrule=0.6pt,
  arc=1.5mm, boxsep=0pt, left=4pt, right=4pt, top=4pt, bottom=4pt, halign=center,
  center, before skip=2pt, after skip=2pt}}
% Biểu tượng "bức ảnh"
\newcommand{\@iconanh}{%
  \begin{tikzpicture}[scale=0.5]
    \draw[bkxanhsang!70, line width=1.2pt, rounded corners=2pt] (0,0) rectangle (2.6,1.9);
    \fill[bkxanhsang!45] (0.25,0.25) -- (1.0,1.15) -- (1.45,0.65) -- (1.8,1.0) -- (2.35,0.25) -- cycle;
    \fill[vang] (1.9,1.45) circle (0.2);
  \end{tikzpicture}}
% Ô chờ ảnh (chưa có tệp) -> tự thay bằng ảnh thật khi tệp tồn tại
\newcommand{\@ochoanh}[1]{%
  \begin{tcolorbox}[enhanced, height=4.4cm, valign=center, halign=center, colback=bkxanhnen,
      frame hidden, borderline={0.7pt}{0pt}{bkxanhsang!70, dash pattern=on 4pt off 2.5pt},
      arc=1mm, boxsep=0pt, left=0pt, right=0pt, top=0pt, bottom=0pt]
    \@iconanh\par\smallskip
    {\small\bfseries\color{bkxanh}KHUNG CHỜ HÌNH ẢNH}\par
    {\footnotesize\color{bkxam}\texttt{figures/\detokenize{#1}}}
  \end{tcolorbox}}
\newcommand{\@timanh}[2]{% #1 độ rộng, #2 tệp
  \IfFileExists{figures/#2}{\includegraphics[width=#1]{#2}}%
   {\IfFileExists{figures/#2.png}{\includegraphics[width=#1]{#2}}%
    {\IfFileExists{figures/#2.jpg}{\includegraphics[width=#1]{#2}}%
     {\IfFileExists{figures/#2.pdf}{\includegraphics[width=#1]{#2}}%
      {\@ochoanh{#2}}}}}}
% \@anh{có khung? true/false}{độ rộng}{tệp}
\newcommand{\@anh}[3]{%
  \ifstrequal{#1}{true}%
    {\begin{tcolorbox}[khunganh, width=#2]\@timanh{\linewidth}{#3}\end{tcolorbox}}%
    {\@timanh{#2}{#3}}}
% \chenhinh[độ rộng]{tệp trong figures/}{Chú thích}{fig:nhan}   -> có khung
% \chenhinh*[...]{...}{...}{...}                                  -> không khung
\NewDocumentCommand{\chenhinh}{s O{0.75\linewidth} m m m}{%
  \begin{figure}[H]\centering
    \IfBooleanTF{#1}{\@anh{false}{#2}{#3}}{\@anh{true}{#2}{#3}}%
    \caption{#4}\label{#5}\end{figure}}
% Hai hình cạnh nhau (có khung):
% \haihinh{tệp1}{chú thích 1}{nhãn1}{tệp2}{chú thích 2}{nhãn2}{chú thích chung}{nhãn chung}
\newcommand{\haihinh}[8]{%
  \begin{figure}[H]\centering
    \begin{subfigure}[t]{0.48\linewidth}\centering\@anh{true}{\linewidth}{#1}\caption{#2}\label{#3}\end{subfigure}\hfill
    \begin{subfigure}[t]{0.48\linewidth}\centering\@anh{true}{\linewidth}{#4}\caption{#5}\label{#6}\end{subfigure}
    \caption{#7}\label{#8}
  \end{figure}}

% ==============================================================================
% 6. MÃ NGUỒN
% ==============================================================================
% Badge ngôn ngữ (góc trên phải khung mã)
\NewDocumentCommand{\@badgema}{O{bkxanhsang} m}{%
  \tcbox[on line, colback=#1, colframe=#1, colupper=white, boxrule=0pt, arc=3pt,
         boxsep=0pt, left=5pt, right=5pt, top=2pt, bottom=2pt,
         fontupper=\scriptsize\bfseries\sffamily]{#2}}
% Biểu tượng tệp
\newcommand{\@icontep}{%
  \tikz[baseline=-0.25ex]{\draw[bkxam, line width=0.7pt, line join=round]
    (0,-0.11) -- (0,0.17) -- (0.13,0.17) -- (0.2,0.10) -- (0.2,-0.11) -- cycle;}}

% 6.1 Khung mã:  \begin{khungma}{tên_tệp.py}{Python}[tùy chọn listings]
%     #1 tên tệp (góc trái)  #2 ngôn ngữ (badge góc phải, cũng là language của listings)
\NewTCBListing{khungma}{m m !O{}}{enhanced, breakable, listing only,
  colback=manen, colframe=bkxam!40, boxrule=0.6pt, arc=2mm,
  left=4pt, right=6pt, top=3pt, bottom=3pt, before skip=8pt, after skip=8pt,
  colbacktitle=bkxam!10, coltitle=bkxanhdam, titlerule=0.5pt, toptitle=2pt, bottomtitle=2pt,
  title={\@icontep\hspace{5pt}\ttfamily\small\bfseries\detokenize{#1}\hfill\@badgema{#2}},
  listing options={style=mauma, frame=none, xleftmargin=18pt, numbersep=9pt,
                   aboveskip=3pt, belowskip=1pt, language={#2}, #3}}

% 6.2 Khung lệnh bash (kiểu cửa sổ terminal):  \begin{khungbash}[Tiêu đề] ... \end{khungbash}
\definecolor{termnen}{HTML}{1E2230}
\definecolor{termthanh}{HTML}{2D3346}
\definecolor{termchu}{HTML}{E6E9F0}
\definecolor{termxanh}{HTML}{7EE787}
\definecolor{termcyan}{HTML}{79C0FF}
\definecolor{termvang}{HTML}{F2CC60}
\lstdefinestyle{terminal}{language=bash, backgroundcolor={},
  basicstyle=\ttfamily\small\color{termchu}, keywordstyle=\color{termcyan}\bfseries,
  commentstyle=\color{termchu!55!termnen}\itshape, stringstyle=\color{termvang},
  numbers=none, frame=none, xleftmargin=0pt, columns=fullflexible,
  literate={\$}{{\textcolor{termxanh}{\$}}}1}
\newcommand{\@chamterm}{%
  \tikz[baseline=-0.6ex]{\fill[do!80!white] (0,0) circle (3.3pt);
    \fill[vang] (0.34,0) circle (3.3pt); \fill[xanhla!70!white] (0.68,0) circle (3.3pt);}}
\NewTCBListing{khungbash}{!O{Terminal}}{enhanced, breakable, listing only,
  colback=termnen, colframe=termthanh, boxrule=0.6pt, arc=2mm,
  left=10pt, right=10pt, top=5pt, bottom=6pt, before skip=8pt, after skip=8pt,
  colbacktitle=termthanh, coltitle=termchu!75, titlerule=0pt, toptitle=3pt, bottomtitle=3pt,
  fonttitle=\ttfamily\footnotesize,
  title={\@chamterm\hfill #1\hfill\@badgema[termxanh!55!black]{bash}},
  listing options={style=mauma, style=terminal, aboveskip=3pt, belowskip=1pt}}

% 6.3 Tự định nghĩa ngôn ngữ listings chưa có sẵn
\lstdefinelanguage{JavaScript}{
  morekeywords={break,case,catch,class,const,continue,default,do,else,export,extends,
    false,finally,for,function,if,import,in,instanceof,let,new,null,return,super,
    switch,this,throw,true,try,typeof,undefined,var,while,async,await,of},
  sensitive=true, morecomment=[l]{//}, morecomment=[s]{/*}{*/},
  morestring=[b]', morestring=[b]", morestring=[b]`}
\lstdefinelanguage{JSON}{morestring=[b]", sensitive=true}

% Ví dụ LaTeX: hiện mã (trên) và kết quả (dưới) — dùng trong Phụ lục B
\NewTCBListing{vidulatex}{!o}{enhanced, breakable, listing and text,
  colback=manen, colbacklower=white, colframe=bkxam!45, boxrule=0.5pt, arc=1.5mm,
  segmentation style={bkxam!45, dashed}, left=8pt, right=8pt,
  IfValueT={#1}{title={#1}, colbacktitle=bkxam!15, coltitle=bkxanh, fonttitle=\bfseries\small},
  listing options={style=mauma, language={[LaTeX]TeX}, numbers=none, frame=none,
    xleftmargin=0pt, basicstyle=\ttfamily\footnotesize, keywordstyle=\color{matukhoa}}}

% ==============================================================================
% 7. THAM CHIẾU NHANH
% ==============================================================================
\newcommand{\Hinh}[1]{\figurename~\ref{#1}}
\newcommand{\Bang}[1]{\tablename~\ref{#1}}
\newcommand{\PT}[1]{phương trình~\eqref{#1}}
\newcommand{\Chuong}[1]{Chương~\ref{#1}}
\newcommand{\Muc}[1]{mục~\ref{#1}}
\newcommand{\PhuLuc}[1]{\appendixname~\ref{#1}}

% ==============================================================================
% 8. DANH MỤC TỪ VIẾT TẮT
% ==============================================================================
% \begin{danhmucviettat} \viettat{AI}{Trí tuệ nhân tạo} ... \end{danhmucviettat}
\newenvironment{danhmucviettat}
  {\rowcolors{2}{white}{bkxanhnen}%
   \begin{longtable}{@{}L{3.2cm}L{\dimexpr\linewidth-3.2cm-2\tabcolsep\relax}@{}}}
  {\end{longtable}}
\newcommand{\viettat}[2]{\textbf{\textcolor{bkxanh}{#1}} & #2 \\}

\makeatother
